%%%%%%%%%%%%%%%%%%%%%%%%%%%%%%%%%%%%%%%%%%%%%%%%%%%%%%%
% % Lines starting with % are comments, which are ignored.
% % This is a handy way of indicating the date and version of
% % your document, to wit:
% %
% % LaTeX sample file
% % Modified March, 2002
% %
%%%%%%%%%%%%%%%%%%%%%%%%%%%%%%%%%%%%%%%%%%%%%%%%%%%%%%%
\DocumentMetadata{
  % PDF standards: ONE key whose value is a comma-separated list --
  % UA-2 (accessibility) plus A-4f (archival; the archival companion of
  % UA-2). Do NOT use separate pdfstandard=... keys: with ordinary
  % key-value parsing only the last one takes effect.
  pdfstandard={UA-2,A-4f},
  lang=en-US,
  tagging=on,
  tagging-setup={
    math/setup=mathml-SE
  }
}

\documentclass[12pt]{exam}
% worksheet-macros.tex -- shared package setup for worksheet-template.tex.
%
% NOTE ON PATHS: every template in this repository is stored one level up
% from where it is compiled (see README.md). All relative \input paths
% below are resolved against the folder the document is compiled in,
% which sits one directory above this file.

% --- Math ------------------------------------------------------------
% amsmath: equation environments and numbering
% amsfonts: \mathbb etc.  |  mathrsfs: \mathfrak  |  stmaryrd: extra symbols
\usepackage{amsmath}
\usepackage{amsfonts}
\usepackage{mathrsfs}
\usepackage{stmaryrd}

% --- Layout and glyphs ------------------------------------------------
% float: [H] float placement (stop floats where you put them)
% fontspec: Unicode/OpenType fonts (compile with XeLaTeX or LuaLaTeX)
% emoji: emoji glyphs in the text
\usepackage{float}
\usepackage{fontspec}
\usepackage{emoji}

% --- Typography --------------------------------------------------------
% Atkinson Hyperlegible as the document-wide sans serif (sfdefault):
% a high-legibility face that is easier to read for students with
% dyslexia or low vision.
\usepackage[sfdefault]{atkinson}

% --- Colors -------------------------------------------------------------
% dvipsnames: extended named color palette (the Accessible* colors in
%   lecture-macros.tex are built on it);
% table: \rowcolors / \cellcolor for shading worksheet tables.
% NOTE: xcolor may only be loaded ONCE per document -- a second
% \usepackage{xcolor} with different options is a fatal "Option clash"
% error, so ALL of its options are listed here.
\usepackage[dvipsnames,table]{xcolor}

% --- Multimedia ---------------------------------------------------------
% animate: embeds audio/video (and legacy Flash) media in the worksheet.
\usepackage{animate}

% --- Bibliography --------------------------------------------------------
% biblatex (biber backend): citations and a reference list for worksheets.
\usepackage[
  backend=biber,
  style=alphabetic,
  url=true,
  doi=true,
  isbn=true,
  backref=true
]{biblatex}

% --- Links ---------------------------------------------------------------
% hyperref makes every \ref/citation/URL clickable in the PDF.
% Link colors must satisfy WCAG 1.4.3 (>= 4.5:1 contrast on white):
%   blue = 8.6:1, green!40!black = 7.2:1.
% (The old cyan/magenta values measured 1.25:1 / 3.1:1 and failed.)
\usepackage{hyperref}
\hypersetup{
  colorlinks=true,
  linkcolor=blue,
  filecolor=blue,
  urlcolor=green!40!black,
}

% --- Shared course macros ------------------------------------------------
% Root macros.tex: math shorthands used across all course documents
% (\RR, \NN, \ZZ, \QQ, \grad, \curl, \tr, \der, \pder, ...).
\newcommand{\RR}{\mathbb{R}}
\newcommand{\NN}{\mathbb{N}}
\newcommand{\ZZ}{\mathbb{Z}}
\newcommand{\QQ}{\mathbb{Q}}
\newcommand{\vv}{\vec{v}}
\newcommand{\ww}{\vec{w}}
\newcommand{\xx}{\vec{x}}
\newcommand{\yy}{\vec{y}}
\newcommand{\zz}{\vec{z}}
\newcommand{\uu}{\vec{u}}
\newcommand{\basis}[1]{\vec{e_{#1}}}
\newcommand{\id}{\mathrm{id}}
\newcommand{\CC}{\mathbb{C}}
\renewcommand{\div}{\operatorname{div}}
\newcommand{\grad}{\operatorname{grad}}
\newcommand{\curl}{\operatorname{curl}}
\newcommand{\tr}{\operatorname{tr}}
\newcommand{\der}[2]{\dfrac{d #1}{d #2}}
\newcommand{\pder}[3][]{\dfrac{\partial^{#1} #2}{\partial #3}}


% % Title and author(s)
%%%%%%%%%%%%%%%%%%%%%%%%%%%%%%%%%%%%%%%%%%%%%%%%%%%%%%%
\title{Worksheet 4 - Limits at Infinity and the Squeeze Theorem}
\author{MATH 1271}
\date{09/17/2026}
%%%%%%%%%%%%%%%%%%%%%%%%%%%%%%%%%%%%%%%%%%%%%%%%%%%%%%%
% %
% % The next command allows your in import encapsulated
% % postscript files, .epsf or .eps files, which
% % contain vector graphic image data.
% %
%%%%%%%%%%%%%%%%%%%%%%%%%%%%%%%%%%%%%%%%%%%%%%%%%%%%%%%
% %   The first thanks indicates your affiliation
% %
% %  Just the name here.
% %
% % Your mailing address goes at the end.
% %
% % \thanks is also how you indicate grant support
% %
%%%%%%%%%%%%%%%%%%%%%%%%%%%%%%%%%%%%%%%%%%%%%%%%%%%%%%%

%%%%%%%%%%%%%%%%%%%%%%%%%%%%%%%%%%%%%%%%%%%%%%%%%%%%%%%
% $ The next command ensures the margins are 1 inch on
% % all sides.
% %
%%%%%%%%%%%%%%%%%%%%%%%%%%%%%%%%%%%%%%%%%%%%%%%%%%%%%%%
\begin{document}
\newpage
\maketitle

\section{Limits Involving Infinity}
\begin{enumerate}

\item[1.] In this problem, we consider the function $f(x)=\dfrac{4-x^2}{x^2+x-6}$.

\begin{enumerate}

\item[(a)] Use algebra to find the value of each of the following limits. If a limit does not exist, write ``DNE" and then explain why it does not exist.

	\begin{enumerate}
	\item[(i)] $\displaystyle\lim_{x\to 2}f(x)=$
	\vfill
	\item[(ii)] $\displaystyle\lim_{x\to -3}f(x)=$
	\vfill
	\item[(iii)] $\displaystyle\lim_{x\to+\infty} f(x)=$
	\vfill
	\item[(iv)] $\displaystyle\lim_{x\to -\infty} f(x)=$
	\vfill
	\end{enumerate}
	
\item[(b)] How many vertical asymptotes does $f(x)$ have? Find the equations of all vertical asymptotes of $f(x)$. 

\vspace{1in}

\item[(c)] How many horizontal asymptotes does $f(x)$ have? Find the equations of all horizontal asymptotes of $f(x)$.

\vspace{1in}

\end{enumerate}

\newpage

\item[2.] Each of the following limits is some sort of infinity (either $+\infty$ or $-\infty$). Determine which type. Discuss your reasoning with your group.

\begin{enumerate}

\item[(a)] $\displaystyle\lim_{x\to 5^+}\frac{x^3}{x-5}$

\vfill

\item[(b)] $\displaystyle\lim_{x\to 5^-}\frac{x^3}{x-5}$

\vfill

\item[(c)] $\displaystyle \lim_{x \rightarrow -\infty} \frac{x^3}{x-5}$
\vfill

\item[(d)] $\displaystyle\lim_{x\to -2^-}\frac{x^2-9}{x^2+5x+6}$

\vfill

\item[(e)] $\displaystyle\lim_{x\to 1^+}\ln(x^2-1)$

\vfill


{\bf Checking understanding:} Discuss the following questions with your group:

	\begin{itemize}	
	\item What if part (d) had asked for $\displaystyle\lim_{x\to -3}\frac{x^2-9}{x^2+5x+6}$ instead? How does that change your method and/or final answer?

  \item What if part (c) had asked for $x \to \infty$ instead? How does that change your method and/or final answer?
	
	\item Does it make sense to ask about the one-sided limit $\displaystyle\lim_{x\to 1^-}\ln(x^2-1)$? Why or why not? (Hint: What is the domain of the function $f(x)=\ln(x^2-1)$?)
	
	\end{itemize}

\end{enumerate}
\newpage
\item[3.] You may assume that the function $y=\dfrac{1-\cos(x)}{x^2}$ satisfies the inequalities
$$\frac{1}{2}-\frac{x^2}{24}<\frac{1-\cos(x)}{x^2}<\frac{1}{2}$$
for all values of $x$ close to $0$. (You can learn where this comes from in Calculus II!) Use this information as needed to find the exact values of each limit below (if a limit does not exist, write ``DNE"). If the Squeeze Theorem is used, show all work; if not, then explain why the theorem does not apply.

\begin{enumerate}

\item[(a)] $\displaystyle\lim_{x\to 0}\frac{1-\cos(x)}{x^2}$

\vfill

\item[(b)] $\displaystyle\lim_{x\to \pi}\frac{1-\cos(x)}{x^2}$

\vfill

\item[(c)] $\displaystyle \lim_{x \to \infty} \frac{1 - \cos(x)}{x^2}$

\vfill

\end{enumerate}

\newpage
\item[4.]

  Suppose we know that the function $f(x) = \dfrac{\sin(x)}{x}$
  satisfies the following inequality when $-\dfrac{\pi}{2} \leq x \leq \dfrac{\pi}{2}$:

  \begin{equation*}
    \cos(x) \leq \dfrac{\sin(x)}{x} \leq 1
  \end{equation*}

  (See Apex Calculus Example 1.3.9 for why this holds). Compute the values of the following limits:

  \begin{enumerate}
    \item[(a)] $\displaystyle \lim_{x\rightarrow 0} \dfrac{\sin(x)}{x}$
      \vfill

    \item[(b)] $\displaystyle \lim_{x \rightarrow 2\pi} \dfrac{\sin(x)}{x}$ (note, the value of this limit is \textit{NOT} $1$).
      \vfill
  \end{enumerate}

\item[5.] Suppose that a particle is oscillating with position described as a function of time by:
  \begin{equation*}
    s(t) = \sin(t)
  \end{equation*}

  Use the previous problem to determine the instantaneous velocity of the particle when $t = 0$.
  \vfill
\end{enumerate}







% Most people don't use these, so they are "commented out"
% by starting the lines with a "%"
%\begin{keywords}
%   \LaTeX, typesetting
%\end{keywords}

%\begin{AMS}
%   50C60, 18C25
%\end{AMS}

%%%%%%%%%%%%%%%%%%%%%%
% % Here is the start of the Text
%%%%%%%%%%%%%%%%%%%%%%
\end{document}
